\section{Performance and Benchmarks}

\subsection{Benchmark Method}

The post-gadget cross-host benchmark used one kernel build, one
benchmark bundle, the same orchestrator, performance governor, and the
same three tiers across eight lab hosts:
\begin{itemize}
  \item \textbf{micro}: \texttt{getpid-loop}, RDTSC-bracketed, pure
        trap-mechanism cost.
  \item \textbf{py}: canned Python startup/mixed-syscall workload.
  \item \textbf{stress}: mt-mini SMP T=8, ncpus=4, strict memory
        verification.
\end{itemize}

\subsection{Micro Tier}

\begin{figure}[h]
\centering
\begin{tikzpicture}
\begin{axis}[
  width=0.92\linewidth,
  height=6cm,
  ybar,
  ymin=0,
  ylabel={speedup over seccomp (x)},
  symbolic x coords={s0,s1,s2,s3,s4,s5,s6,s7},
  xtick=data,
  bar width=13pt,
  fill=UMLGreen!70,
  draw=UMLGreen,
  nodes near coords,
  nodes near coords style={font=\scriptsize,rotate=90,anchor=west},
  grid=major,
  major grid style={draw=UMLGrid},
]
\addplot table[col sep=semicolon,x=host,y=micro]{../data/bench_speedups.csv};
\end{axis}
\end{tikzpicture}
\caption{Post-gadget micro speedup. Trivial syscall dispatch is 193--908x faster than seccomp across the eight hosts.}
\end{figure}

The micro result is the clearest evidence that the LSTAR gadget
changed the backend's cost class. Before the gadget, a trivial syscall
paid a host-exit round trip. After gadget revival, the measured
\KVMvTwo{} cost is roughly 87--153 cycles on the tested machines,
while seccomp remains tens of thousands of cycles per call.

\subsection{Application and Stress Tiers}

\begin{figure}[h]
\centering
\begin{tikzpicture}
\begin{axis}[
  width=0.92\linewidth,
  height=6cm,
  ybar,
  ymin=0,
  ymax=10,
  ylabel={speedup over seccomp (x)},
  symbolic x coords={s0,s1,s2,s3,s4,s5,s6,s7},
  xtick=data,
  bar width=7pt,
  legend style={at={(0.5,1.05)},anchor=south,legend columns=2,draw=none},
  grid=major,
  major grid style={draw=UMLGrid},
]
\addplot+[fill=UMLBlue!65,draw=UMLBlue] table[col sep=semicolon,x=host,y=py]{../data/bench_speedups.csv};
\addplot+[fill=UMLAmber!75,draw=UMLAmber] table[col sep=semicolon,x=host,y=stress]{../data/bench_speedups.csv};
\legend{Python tier,stress tier}
\end{axis}
\end{tikzpicture}
\caption{Application and stress tiers. \KVMvTwo{} wins on every host in both tiers.}
\end{figure}

The Python tier matters because it turns the micro win into an
application-level win: 2.99--4.06x faster than seccomp. The stress
tier matters because it combines performance with correctness: the
run reported zero strict memory failures and zero verify failures on
every host.

\subsection{Detailed Cross-Host Table}

\begin{table}[h]
\centering
\caption{Post-gadget speedup over seccomp.}
\small
\begin{tabular}{@{}llrrr@{}}
\toprule
Host & CPU & micro & Python & stress \\
\midrule
s0 & i9-12900K & 284x & 3.73x & 7.26x \\
s1 & Xeon E3-1225 v6 & 743x & 3.05x & 3.71x \\
s2 & Xeon E3-1225 v5 & 813x & 2.99x & 3.94x \\
s3 & Xeon W-2123 & 707x & 3.14x & 4.16x \\
s4 & i5-12600K & 193x & 2.99x & 8.72x \\
s5 & Ryzen 7 7840HS & 861x & 3.87x & 4.30x \\
s6 & Ryzen 7 7840HS & 908x & 4.06x & 4.32x \\
s7 & Ryzen 7 7840HS & 904x & 3.91x & 4.89x \\
\bottomrule
\end{tabular}
\end{table}

\subsection{Performance Interpretation}

The performance story has three layers:
\begin{itemize}
  \item \textbf{Mechanism win.} The micro tier proves that
        in-guest fast-path dispatch removes most seccomp signal/stub
        cost for trivial syscalls.
  \item \textbf{Workload win.} The Python tier proves that real
        userland startup benefits from those trivial syscalls becoming
        cheap.
  \item \textbf{Concurrency win.} The stress tier is mmap/page-fault
        heavy, so the gadget is less directly relevant; the continuing
        3.71--8.72x speedup comes from the broader vCPU and SMP design.
\end{itemize}

The correct benchmark unit is the ratio against seccomp on the same
host, not the raw absolute latency across machines. That choice
removes most CPU-frequency and microarchitecture noise and keeps the
claim tied to a backend decision.
